%======================================================================
\chapter{Two descriptions give a homeomorphism}\label{ch:two}
%======================================================================

This chapter proves Part~I of the uniqueness theorem.

\begin{theorem}[Bing 1948, Theorem~12]\label{thm:two-desc}
Let $X_1,X_2$ be compact metric spaces, and for $c=1,2$ let $D^{(c,0)},D^{(c,1)},\dots$ be a description from $P_c$ to $Q_c$ in
$X_c$, with $P_c\neq Q_c$ and described set $M_c$. Then there is a homeomorphism $T\colon M_1\to M_2$ with $T(P_1)=P_2$ and
$T(Q_1)=Q_2$.
\end{theorem}

\formal{\lean{HPC.theorem12}, which combines \lean{HPC.Description.exists_homeomorph} (the theorem assuming pattern following)
with \lean{HPC.Description.patternFollowing}.}

The idea: build, alternately in the two spaces, chains $C^{(1,i)}$ in $X_1$ and $C^{(2,i)}$ in $X_2$ with the same number of links
$N_i$, such that $C^{(1,i+1)}$ and $C^{(2,i+1)}$ follow the \emph{same} pattern in $C^{(1,i)}$ and $C^{(2,i)}$. Then the links of
these chains are matching covers (Chapter~\ref{ch:matching}). At each step one side is refined by a deep level of its own
description, which yields a pattern, and the other side follows this pattern by Theorem~\ref{thm:pattern-following}.

\section{Patterns from refinements, and matching}

\begin{lemma}\label{lem:pattern-of-refines}
Let $Y$ (of length $N$) and $Z$ (of length $m$) be chains from $P$ to $Q$, with $Y$ refining $Z$. Then any refinement map
$x\colon[N]\to[m]$ is a pattern in $Z$, and $Y$ follows $x$ in $Z$.
\end{lemma}

\begin{proof}
$x$ has unit steps by \vone{Lemma~2.4}. $P\in Y_0\subseteq Z_{x(0)}$ and $P$ lies only in $Z_0$, so $x(0)=0$; likewise
$x(N-1)=m-1$ using $Q$. The last statement is the definition of a refinement map.
\end{proof}

\formal{\lean{HPC.Chain.exists_pattern_of_refines}.}

\begin{lemma}[Matching]\label{lem:pattern-matching}
Let $Y_1,Z_1$ be chains in $X_1$ and $Y_2,Z_2$ chains in $X_2$, with $Y_1,Y_2$ of length $N$ and $Z_1,Z_2$ of length $m$. Suppose $Y_1$
follows a function $x\colon[N]\to[m]$ in $Z_1$ and $Y_2$ follows the same $x$ in $Z_2$, and the links of $Z_2$ are bounded with
diameter $<\varepsilon$. If $(Y_1)_j$ meets $(Z_1)_k$, then some point of $(Y_2)_j$ is at distance $<\varepsilon$ from some point of
$(Z_2)_k$.
\end{lemma}

\begin{proof}
A point of $(Y_1)_j\cap(Z_1)_k$ lies in $(Z_1)_{x(j)}\cap(Z_1)_k$, so $\abs{x(j)-k}\le1$, and therefore $(Z_2)_{x(j)}$ meets $(Z_2)_k$,
say at $q$. Take any $p\in(Y_2)_j$. Then $p,q\in(Z_2)_{x(j)}$, so $d(p,q)\le\diam(Z_2)_{x(j)}<\varepsilon$.
\end{proof}

\formal{\lean{HPC.Chain.pattern_matching}.}

\begin{lemma}\label{lem:consolidation-meets}
Let $C$ be a consolidation of a level $D^{(l)}$ of a description in a compact space, with described set $M$. Then
$M\subseteq\bigcup C$, and every link of $C$ meets $M$.
\end{lemma}

\begin{proof}
$M\subseteq\bigcup D^{(l)}\subseteq\bigcup C$ since $D^{(l)}$ refines $C$. A link of $C$ is a nonempty union of links of $D^{(l)}$, and
every link of $D^{(l)}$ meets $M$ by \vone{Proposition~3.9}.
\end{proof}

\formal{\lean{HPC.Description.M_subset_consolidation}, \lean{HPC.Description.consolidation_link_inter_M_nonempty},
\lean{HPC.Description.link_inter_M_nonempty}.}

\begin{lemma}\label{lem:level-refining}
Let $C$ be a consolidation of the level $D^{(l)}$ and $\delta>0$. There is $k\ge l$ such that $D^{(k)}$ refines $C$ and
$\varepsilon_k\le\delta$.
\end{lemma}

\begin{proof}
Choose $k\ge l$ with $\varepsilon_k\le\delta$. Then $D^{(k)}$ refines $D^{(l)}$ (\vone{Lemma~3.2(b)}), which refines $C$.
\end{proof}

\formal{\lean{HPC.Description.exists_level_refining}.}

\section{The alternating construction}

Fix the two descriptions of Theorem~\ref{thm:two-desc}; write $n^c_l$ and $\varepsilon^c_l$ for the lengths and mesh bounds
of the description in $X_c$, and $D^{(c,l)}$ for its levels. By Theorem~\ref{thm:pattern-following} both descriptions have the
pattern following property.

\begin{definition}\label{def:rung}
A \emph{rung} consists of a number $N$, chains $C^1$ in $X_1$ and $C^2$ in $X_2$, both of length $N$, and numbers $m>0$,
$m_f>0$, such that
\begin{enumerate}[label=(R\arabic*)]
\item $C^c$ goes from $P_c$ to $Q_c$ and is a consolidation of some level of the description in $X_c$ ($c=1,2$);
\item every link of $C^1$ and of $C^2$ has diameter $\le m$;
\item one of the two chains, the \emph{active} one, is itself a level of its description, and its links have diameter $\le m_f$.
\end{enumerate}
\end{definition}

\formal{\lean{HPC.Rung}. Formally, (R3) asks that the active side be able to follow patterns
(\lean{HPC.Description.CanFollow}), which holds for every level by pattern following
(\lean{HPC.Description.canFollow_level}).}

\begin{lemma}[The first rung]\label{lem:base-rung}
There is a rung with $m_f=1/2$.
\end{lemma}

\begin{proof}
Choose a level $a$ of the first description with $\varepsilon^1_a<1/2$ and $\varepsilon^1_a<d(P_1,Q_1)$. Then $n^1_a\ge2$, since $P_1$
and $Q_1$ cannot lie in one link. By Lemma~\ref{lem:n-tendsto} choose $b$ with $n^2_b\ge n^1_a$, and by
Lemma~\ref{lem:consolidate-length} a consolidation $F$ of $D^{(2,b)}$ of length $n^1_a$ from $P_2$ to $Q_2$. Take $C^1=D^{(1,a)}$
(active, $m_f=1/2$), $C^2=F$, and $m=\max(1/2,\diam X_2+1)$; $X_2$ is compact, hence bounded.
\end{proof}

\formal{\lean{HPC.Rung.exists_base}.}

\begin{lemma}[The next rung]\label{lem:next-rung}
Given a rung $(N,C^1,C^2,m,m_f)$ and $\delta>0$, there are a rung $(N',C'^1,C'^2,m',\delta)$ with $m'=\max(\delta,m_f)$ and a
function $x\colon[N']\to[N]$ such that $C'^1$ follows $x$ in $C^1$ and $C'^2$ follows $x$ in $C^2$.
\end{lemma}

\begin{proof}
Suppose $C^1$ is active, $C^1=D^{(1,a)}$ (the other case is symmetric). The inactive side $C^2$ is a consolidation of some level of
the second description; by Lemma~\ref{lem:level-refining} there is a level $c$ such that $D^{(2,c)}$ refines $C^2$ and
$\varepsilon^2_c\le\delta$. By Lemma~\ref{lem:pattern-of-refines} a refinement map $x\colon[n^2_c]\to[N]$ is a pattern in $C^2$, followed
by $D^{(2,c)}$. Since $C^1=D^{(1,a)}$ is a level, pattern following (Theorem~\ref{thm:pattern-following}) gives $a'>a$ and a
consolidation $C'^1$ of $D^{(1,a')}$ of length $n^2_c$, from $P_1$ to $Q_1$, following $x$ in $C^1$. Let $N'=n^2_c$ and
$C'^2=D^{(2,c)}$, now the active side, with $m_f'=\delta$. The links of $C'^2$ have diameter $\le\varepsilon^2_c\le\delta$; each link of
$C'^1$ lies in a link of $C^1$, of diameter $\le m_f$. So $m'=\max(\delta,m_f)$ bounds all diameters.
\end{proof}

\formal{\lean{HPC.Rung.exists_next}, via \lean{HPC.Description.exists_half_step}.}

\section{Proof of Theorem~\ref{thm:two-desc}}

\begin{proof}
Start with the rung $R_0$ of Lemma~\ref{lem:base-rung} and let $R_{k+1}$ be obtained from $R_k$ by Lemma~\ref{lem:next-rung}
with $\delta=2^{-(k+2)}$. Write $R_i=(N_i,C^{(1,i)},C^{(2,i)},m_i,m_{f,i})$. Then $m_{f,i}=2^{-(i+1)}$ for all $i$, and for
$i\ge1$, $m_i=\max(2^{-(i+1)},2^{-i})=2^{-i}$. Put $\varepsilon_i=2m_i$; these numbers are positive and summable.

We check that the links $x^c_{i,k}=C^{(c,i)}_k$ ($k\in[N_i]$) are matching covers of $M_1$ and $M_2$ for $\varepsilon$
(Definition~\ref{def:matching}). The links are open; \ref{M:cover} and~\ref{M:meet} hold by Lemma~\ref{lem:consolidation-meets};
\ref{M:diam} holds since diameters are at most $m_i<\varepsilon_i$ (links are bounded because $X_c$ is compact). For~\ref{M:match}, by
Lemma~\ref{lem:next-rung} there is $x$ such that $C^{(1,i+1)}$ and $C^{(2,i+1)}$ follow $x$ in $C^{(1,i)}$ and $C^{(2,i)}$
respectively; Lemma~\ref{lem:pattern-matching} gives~\ref{M:match} in both directions, with $\varepsilon=\varepsilon_i$.

Also $M_1$ and $M_2$ are compact (closed subsets of compact spaces, \vone{Lemma~3.2(d)}). Theorem~\ref{thm:matching} gives a
homeomorphism $T\colon M_1\to M_2$ with the tracking property. The constant sequence $a_i=0$ tracks $P_1$ and $P_2$, because each
$C^{(c,i)}$ goes from $P_c$ to $Q_c$; so $T(P_1)=P_2$. Likewise the sequence $a_i=N_i-1$ tracks $Q_1$ and $Q_2$, so $T(Q_1)=Q_2$.
\end{proof}

\section{Descriptions in non-compact spaces}

Theorem~\ref{thm:two-desc} assumes that the ambient spaces are compact. It is enough that the links lie in compact sets.

\begin{corollary}\label{cor:two-desc-K}
Let $D^{(c,i)}$ be descriptions from $P_c$ to $Q_c\ne P_c$ in arbitrary metric spaces $X_c$ ($c=1,2$), and suppose all links of the
$c$-th description lie in a compact set $K_c\subseteq X_c$. Then there is a homeomorphism $T\colon M_1\to M_2$ with $T(P_1)=P_2$ and
$T(Q_1)=Q_2$.
\end{corollary}

\begin{proof}
Regard each link $D^{(c,i)}_t\subseteq K_c$ as a subset of the compact metric space $K_c$. It is open in $K_c$, the chain conditions
and the conditions ``from $P_c$ to $Q_c$'', ``crooked'' and the diameter bounds do not change, and the closure of a link in $K_c$ is its
closure in $X_c$ intersected with $K_c$, so closures of links are still nested. Thus we have descriptions in the compact spaces $K_c$,
with the same described sets. Apply Theorem~\ref{thm:two-desc}.
\end{proof}

\formal{\lean{HPC.Plane.theorem12_of}, using the restriction \lean{HPC.Description.restrict} (in
\lean{HPC/Plane/Transfer.lean}).}

\begin{remark}
Bing states Theorem~12 under the assumption that the described sets are compact. This is not sufficient: \vone{Remark~3.5} gives a
description, in a non-complete space, whose described set is compact but not connected; it is not homeomorphic to the
pseudo-arc described in the plane. The formal development records this as \lean{HPC.Plane.theorem12_of_false}.
\end{remark}

\section*{Exercises}
\addcontentsline{toc}{section}{Exercises}

\begin{exercise}
Show that in Theorem~\ref{thm:two-desc} one may take the two descriptions to be the same. What does the theorem then say?
(Compare Chapter~\ref{ch:main}.)
\end{exercise}

\begin{exercise}
Why does the construction alternate between the two sides? Show that if one only ever refines the first side, the
mesh of the second side need not tend to zero.
\end{exercise}

\begin{exercise}
Show that the homeomorphism $T$ of Theorem~\ref{thm:two-desc} can be chosen so that also $T(R_1)=R_2$ for given points $R_c$ that
lie, for every $i$, in the $k_i$-th link of $C^{(c,i)}$ for the same indices $k_i$. (This is the tracking clause.)
\end{exercise}
